% Generated draft: requires mathematical and source review.
\documentclass[11pt]{article}
\usepackage[margin=25mm]{geometry}
\usepackage{fontspec}
\setmainfont{Latin Modern Roman}
\usepackage{amsmath,amssymb,mathtools,unicode-math}
\setmathfont{Latin Modern Math}
\usepackage{enumitem,needspace,xurl,hyperref}
\setlength{\parindent}{0pt}
\setlength{\parskip}{5pt}
\setlength{\emergencystretch}{3em}
\newcommand{\N}{\mathbb{N}}
\newcommand{\Z}{\mathbb{Z}}
\newcommand{\Q}{\mathbb{Q}}
\newcommand{\R}{\mathbb{R}}
\newcommand{\C}{\mathbb{C}}
\newcommand{\F}{\mathbb{F}}
\newcommand{\A}{\mathbb{A}}
\newcommand{\PP}{\mathbb{P}}
\newcommand{\E}{\mathbb{E}}
\newcommand{\eps}{\varepsilon}
\newcommand{\dd}{\,\mathrm{d}}
\begin{document}
\section*{UnsolvedMath ID: 700993}
\textbf{Source problem code:} catalogue entry
\subsection{Definitions and mathematical statement}
Let \(\operatorname{Homeo}_{+}(\R)\) be the group, under composition, of strictly increasing homeomorphisms \(\R\to\R\). For \(f,g\in\operatorname{Homeo}_{+}(\R)\), define
\[
f\sim g\quad\Longleftrightarrow\quad \exists T\in\R\ \forall x>T,\quad f(x)=g(x).
\]
The group of germs at \(+\infty\) is
\[
G=\operatorname{Homeo}_{+}(\R)/{\sim},\qquad [f][g]=[f\circ g].
\]
Write \(t(x)=x+1\) and \(\alpha=[t]\).

\textbf{Problem (P. de la Harpe).} Characterize all \(\beta\in G\) such that \(\langle\alpha,\beta\rangle\) is a free group of rank \(2\). Equivalently, characterize those \(\beta=[g]\) for which the homomorphism from the free group \(F(a,b)\) to \(G\), given by \(a\mapsto\alpha\) and \(b\mapsto\beta\), is injective. In terms of representatives, the required condition is that for every nonempty reduced word \(w\) in \(a^{\pm1},b^{\pm1}\),
\[
\forall T\in\R\ \exists x>T,\qquad w(t,g)(x)\ne x.
\]
Here \(w(t,g)\) denotes composition of the indicated homeomorphisms and their inverses. The task is to characterize such germs; this word condition specifies the target property.

\textbf{Supplied background and its qualifications.} The record states that if \(\beta\) is represented by \(x\mapsto x^k\), with \(k\ge3\) an odd integer, then \(\langle\alpha,\beta\rangle\) is free of rank \(2\), and that the cited proofs use Galois theory. It attributes the odd-prime case to S. White and the general odd case \(k\ge3\) to S. A. Adeleke, A. M. W. Glass, and L. Morley. These are supplied background assertions, not independently verified here.

The record also mentions results for every odd \(k\ne\pm1\) and every positive even \(k\), attributed to several papers by S. D. Cohen and A. M. W. Glass. This broader qualification requires clarification: a positive even power is not a homeomorphism of all of \(\R\), although its restriction to a sufficiently positive tail can be extended to an increasing homeomorphism of \(\R\), defining a germ in \(G\). A negative odd power is decreasing on the positive half-line and therefore cannot represent an element of this orientation-preserving germ group. The formulation or ambient setting intended for those negative exponents is not supplied.
\subsection{Short English statement}
Characterize the germs at \(+\infty\) of increasing homeomorphisms of \(\R\) which, together with the germ of translation by \(1\), generate a free group of rank \(2\). Equivalently, no nonempty reduced word in the two generators may equal the identity on an entire sufficiently positive tail. The supplied record labels the problem open and reports odd positive power examples; its wider exponent claims need source clarification.
\subsection{Sources}
\begin{itemize}
\item Supplied local UnsolvedMath registry record: permanent ID \(700993\), exact source code \texttt{catalogue entry}, title \emph{Kourovka Notebook Problem 14.8}; upstream revision \texttt{d371ee917ac9647fbb06cd6fc529deec6f3dda82}. The record supplies the status \texttt{open}; this is a registry label, not an independently verified literature assessment.
\item \emph{The Kourovka Notebook (Unsolved Problems in Group Theory)}, Problem 14.8, proposed by P. de la Harpe, 14th issue (1999); the record identifies the 21st issue edition (2026), \texttt{arXiv:1401.0300v48}. The supplied collection description names E. I. Khukhro and V. D. Mazurov as editors.
\item S. White, \emph{J. Algebra} \textbf{118} (1988), 408--422. Supplied attribution: freeness for odd prime exponents, with a proof relying on Galois theory. No article title or link is supplied.
\item S. A. Adeleke, A. M. W. Glass, and L. Morley, \emph{J. London Math. Soc.} \textbf{43} (1991), 255--268. Supplied attribution: freeness for every odd integer \(k\ge3\), with proofs described as relying on Galois theory. No article title or link is supplied.
\item Several papers by S. D. Cohen and A. M. W. Glass are mentioned for odd \(k\ne\pm1\) and positive even \(k\). No individual titles, publication details, or links are supplied, and the interpretation of the negative-exponent cases in relation to \(G\) remains unresolved.
\end{itemize}
\end{document}
